\documentclass[12pt, a4paper]{ctexart}
\usepackage{amsmath, amssymb, amsthm}
\usepackage{fancyhdr}
\usepackage{geometry}
\usepackage{enumitem}
\usepackage{graphicx}
\usepackage{tikz}
\usepackage{pgfplots}
\pgfplotsset{compat=1.18}

\geometry{left=2cm, right=2cm, top=2.5cm, bottom=2.5cm}
\setCJKmainfont{Hiragino Sans GB}
\setCJKsansfont{Heiti SC}

\pagestyle{fancy}
\fancyhf{}
\fancyhead[L]{培优专题·参数方程与极坐标}
\fancyhead[R]{测试卷}
\fancyfoot[C]{\thepage}
\renewcommand{\headrulewidth}{0.4pt}

\newcommand{\score}[1]{\hfill\makebox[2cm][l]{(#1分)}}

\begin{document}

\begin{center}
\Large\textbf{培优专题测试卷·参数方程与极坐标}

\normalsize
考试时间：120分钟\quad 满分：150分
\end{center}

\vspace{0.5cm}

\noindent\textbf{一、单选题（本大题共8小题，每小题5分，共40分。在每小题给出的四个选项中，只有一项是符合题目要求的。）}

\vspace{0.3cm}

\begin{enumerate}
\item 曲线 $\begin{cases} x = 2\cos\theta \\ y = 3\sin\theta \end{cases}$（$\theta$ 为参数）的普通方程为\score{5}
\begin{enumerate}[label=\Alph*.]
\item $\dfrac{x^2}{4} + \dfrac{y^2}{9} = 1$
\item $\dfrac{x^2}{2} + \dfrac{y^2}{3} = 1$
\item $x^2 + y^2 = 13$
\item $\dfrac{x^2}{9} + \dfrac{y^2}{4} = 1$
\end{enumerate}

\item 极坐标方程 $\rho = 4\cos\theta$ 表示的曲线为\score{5}
\begin{enumerate}[label=\Alph*.]
\item 圆心在 $(2, 0)$，半径为 $2$ 的圆
\item 圆心在 $(0, 2)$，半径为 $2$ 的圆
\item 圆心在 $(4, 0)$，半径为 $4$ 的圆
\item 过原点的直线
\end{enumerate}

\item 参数方程 $\begin{cases} x = t + \dfrac{1}{t} \\ y = t - \dfrac{1}{t} \end{cases}$（$t$ 为参数）表示的曲线为\score{5}
\begin{enumerate}[label=\Alph*.]
\item 椭圆
\item 双曲线
\item 抛物线
\item 圆
\end{enumerate}

\item 将极坐标方程 $\rho\cos\theta = 2$ 化为直角坐标方程为\score{5}
\begin{enumerate}[label=\Alph*.]
\item $x = 2$
\item $y = 2$
\item $x^2 + y^2 = 2$
\item $x + y = 2$
\end{enumerate}

\item 已知直线 $l$ 的参数方程为 $\begin{cases} x = 1 + t\cos\alpha \\ y = 2 + t\sin\alpha \end{cases}$（$t$ 为参数），则 $l$ 过定点\score{5}
\begin{enumerate}[label=\Alph*.]
\item $(1, 2)$
\item $(\cos\alpha, \sin\alpha)$
\item $(2, 1)$
\item $(0, 0)$
\end{enumerate}

\item 极坐标方程 $\rho = \dfrac{2}{1 - \cos\theta}$ 表示的曲线为\score{5}
\begin{enumerate}[label=\Alph*.]
\item 椭圆
\item 双曲线
\item 抛物线
\item 圆
\end{enumerate}

\item 曲线 $\begin{cases} x = 1 + \sqrt{2}\cos\theta \\ y = -1 + \sqrt{2}\sin\theta \end{cases}$ 上的点到直线 $x + y = 0$ 的最大距离为\score{5}
\begin{enumerate}[label=\Alph*.]
\item $\sqrt{2}$
\item $2\sqrt{2}$
\item $2$
\item $3\sqrt{2}$
\end{enumerate}

\item 设 $A(\rho_1, \theta_1)$，$B(\rho_2, \theta_2)$ 为极坐标系中两点，则 $|AB|^2 = $\score{5}
\begin{enumerate}[label=\Alph*.]
\item $\rho_1^2 + \rho_2^2 - 2\rho_1\rho_2\cos(\theta_1 - \theta_2)$
\item $\rho_1^2 + \rho_2^2 + 2\rho_1\rho_2\cos(\theta_1 - \theta_2)$
\item $(\rho_1 - \rho_2)^2 + (\theta_1 - \theta_2)^2$
\item $\rho_1^2 + \rho_2^2 - 2\rho_1\rho_2\sin(\theta_1 - \theta_2)$
\end{enumerate}

\end{enumerate}

\vspace{0.5cm}

\noindent\textbf{二、多选题（本大题共4小题，每小题5分，共20分。在每小题给出的选项中，有多项符合题目要求。全部选对得5分，部分选对得2分，有选错的得0分。）}

\vspace{0.3cm}

\begin{enumerate}[resume]
\item 关于参数方程 $\begin{cases} x = a\cos\theta \\ y = b\sin\theta \end{cases}$（$a > b > 0$），下列正确的有\score{5}
\begin{enumerate}[label=\Alph*.]
\item 表示椭圆 $\dfrac{x^2}{a^2} + \dfrac{y^2}{b^2} = 1$
\item 参数 $\theta$ 是椭圆上点的离心角
\item 当 $\theta = 0$ 时，对应点为 $(a, 0)$
\item 当 $\theta = \dfrac{\pi}{2}$ 时，对应点为 $(0, b)$
\end{enumerate}

\item 关于极坐标，下列正确的有\score{5}
\begin{enumerate}[label=\Alph*.]
\item $\rho = 2a\sin\theta$（$a > 0$）表示圆心在 $(0, a)$、半径为 $a$ 的圆
\item $\theta = \dfrac{\pi}{4}$ 表示过原点、倾角为 $\dfrac{\pi}{4}$ 的直线
\item $\rho = a$（$a > 0$）表示圆心在原点、半径为 $a$ 的圆
\item $\rho = \dfrac{ep}{1 - e\cos\theta}$ 表示以焦点为极点的圆锥曲线
\end{enumerate}

\item 已知直线 $l$ 的极坐标方程为 $\rho\cos\left(\theta - \dfrac{\pi}{4}\right) = \sqrt{2}$，则\score{5}
\begin{enumerate}[label=\Alph*.]
\item $l$ 的直角坐标方程为 $x + y = 2$
\item $l$ 过点 $(1, 1)$
\item $l$ 到原点的距离为 $\sqrt{2}$
\item $l$ 与 $x$ 轴的交点为 $(2, 0)$
\end{enumerate}

\item 设直线 $\begin{cases} x = 1 + t \\ y = 2 - t \end{cases}$ 与圆 $\begin{cases} x = \cos\theta \\ y = \sin\theta \end{cases}$ 交于 $A, B$ 两点，则\score{5}
\begin{enumerate}[label=\Alph*.]
\item 直线的普通方程为 $x + y = 3$
\item 圆的普通方程为 $x^2 + y^2 = 1$
\item $|AB| = \sqrt{2}$
\item $AB$ 中点到原点的距离为 $\dfrac{3}{2}$
\end{enumerate}

\end{enumerate}

\vspace{0.5cm}

\noindent\textbf{三、填空题（本大题共4小题，每小题5分，共20分。）}

\vspace{0.3cm}

\begin{enumerate}[resume]
\item 极坐标方程 $\rho = 2\cos\theta + 2\sqrt{3}\sin\theta$ 化为直角坐标方程为 \underline{\hspace{4cm}}，表示 \underline{\hspace{2cm}}。\score{5}

\item 曲线 $\begin{cases} x = t^2 \\ y = 2t \end{cases}$（$t$ 为参数）的普通方程为 \underline{\hspace{3cm}}。\score{5}

\item 直线 $\begin{cases} x = 2 + t \\ y = -1 + 2t \end{cases}$ 被抛物线 $y^2 = 4x$ 截得的弦长为 \underline{\hspace{2cm}}。\score{5}

\item 极坐标系中，点 $A\left(2, \dfrac{\pi}{3}\right)$ 到直线 $\rho\cos\theta = 3$ 的距离为 \underline{\hspace{2cm}}。\score{5}

\end{enumerate}

\vspace{0.5cm}

\noindent\textbf{四、解答题（本大题共6小题，共70分。解答应写出文字说明、证明过程或演算步骤。）}

\vspace{0.3cm}

\begin{enumerate}[resume]
\item （10分）已知直线 $l$ 的参数方程为 $\begin{cases} x = 1 + t \\ y = -1 + t \end{cases}$（$t$ 为参数），曲线 $C$ 的极坐标方程为 $\rho = 2\cos\theta$。
\begin{enumerate}[label=(\arabic*)]
\item 将 $l$ 和 $C$ 的方程化为普通方程；
\item 求 $l$ 与 $C$ 的交点坐标。
\end{enumerate}

\item （12分）已知椭圆 $C$ 的参数方程为 $\begin{cases} x = 2\cos\theta \\ y = \sin\theta \end{cases}$（$\theta$ 为参数），直线 $l$ 的极坐标方程为 $\rho\cos\left(\theta + \dfrac{\pi}{3}\right) = 1$。
\begin{enumerate}[label=(\arabic*)]
\item 求椭圆 $C$ 的普通方程和离心率；
\item 将直线 $l$ 的方程化为直角坐标方程；
\item 求椭圆 $C$ 上的点到直线 $l$ 的最大距离和最小距离。
\end{enumerate}

\item （12分）已知曲线 $C_1$ 的参数方程为 $\begin{cases} x = 1 + \cos\alpha \\ y = \sin\alpha \end{cases}$（$\alpha$ 为参数），曲线 $C_2$ 的极坐标方程为 $\rho = 4\sin\theta$。
\begin{enumerate}[label=(\arabic*)]
\item 分别将 $C_1$、$C_2$ 化为直角坐标方程；
\item 求 $C_1$ 与 $C_2$ 的交点个数；
\item 求 $C_1$ 上的点与 $C_2$ 上的点之间的最大距离。
\end{enumerate}

\item （12分）在极坐标系中，已知圆 $C: \rho = 2\cos\theta$，直线 $l: \rho\cos\theta - \rho\sin\theta = 3$。
\begin{enumerate}[label=(\arabic*)]
\item 求圆 $C$ 的圆心和半径；
\item 求圆 $C$ 上的点到直线 $l$ 的距离的取值范围；
\item 设 $P$ 为圆 $C$ 上的动点，求 $P$ 到直线 $l$ 的距离的最大值。
\end{enumerate}

\item （12分）已知抛物线 $C: y^2 = 4x$，直线 $l$ 过点 $M(4, 0)$，参数方程为 $\begin{cases} x = 4 + t\cos\alpha \\ y = t\sin\alpha \end{cases}$（$t$ 为参数，$\alpha$ 为倾斜角）。
\begin{enumerate}[label=(\arabic*)]
\item 当 $\alpha = \dfrac{\pi}{4}$ 时，求 $l$ 与 $C$ 的交点坐标及 $|MA| \cdot |MB|$；
\item 证明：$|MA| \cdot |MB|$ 与 $\alpha$ 无关；
\item 求 $\dfrac{1}{|MA|} + \dfrac{1}{|MB|}$ 的取值范围。
\end{enumerate}

\item （12分）已知椭圆 $\dfrac{x^2}{4} + y^2 = 1$，$A$ 为椭圆上一点，参数坐标为 $A(2\cos\theta, \sin\theta)$。
\begin{enumerate}[label=(\arabic*)]
\item 求 $|OA|$（$O$ 为原点）的最大值和最小值；
\item 求椭圆上的点到直线 $x + 2y - 4 = 0$ 的最大距离和最小距离；
\item 设 $B$ 为椭圆上另一点，且 $OA \perp OB$，求 $\triangle OAB$ 面积的最大值。
\end{enumerate}

\end{enumerate}

\end{document}
